\documentclass[nobib]{tufte-book}

\input{castelpre.tex}
%\input{symbols.tex}
\usepackage{pdfpages}

\usepackage{lipsum}
\usepackage{parskip}
\usepackage{titletoc}

\newcommand\circled[1]{
    \begin{tikzpicture}[baseline=(char.base)]%
        \node[circle,draw,inner sep=1pt] (char) {\textsf{#1}};%
\end{tikzpicture}}
% minicircle for in figures!
\newcommand\mc[1]{\footnotesize\circled{#1}}

\usepackage{cmbright}
\usepackage{bm}

% \usepackage{eso-pic}                % put things into background 
% \usepackage{lipsum}                 % for sample text

% \definecolor{reallylightgray}{HTML}{FAFAFA}
% \AddToShipoutPicture{% from package eso-pic: put something to the background
%     \ifthenelse{\isodd{\thepage}}{
%           % ODD page: left bar
%           \AtPageLowerLeft{% start the bar at the left bottom of the page
%             \put(\LenToUnit{\dimexpr\paperwidth-222pt},0){% move it to the top right
%                 \color{reallylightgray}\rule{222pt}{297mm}%
%               }%
%           }%
%       }%
%       {%
%         \AtPageLowerLeft{% put it at the left bottom of the page
%           \color{reallylightgray}\rule{222pt}{297mm}%
%         }%
%    }%
% }

\begin{document}
 \setcounter{page}{0}
\setcounter{chapter}{0}
\chapter{Druck und Strömungen}

\section*{Aufgabe 01}
%\marginpar[]{\color{uniblau} Fluid: \par Gas oder Flüssigkeit}
Ein Fluid durchströmt ein Rohr mit drei unterschiedlichen Querschnittsflächen. 
Im ersten Bereich beträgt die Querschnittsfläche 20 cm$^2$ und die Fließgeschwindigkeit des Fluids beträgt 3 cm/s.
%\marginpar[]{\color{uniblau} $A_1 \cdot v_1 = A_2 \cdot v_2$}
\begin{enumerate}
\item Im zweiten Bereich beträgt die Querschnittsfläche 10 cm$^2$. Wie hoch ist die Geschwindigkeit des Fluids hier?
\item Im dritten Bereich beträgt die Querschnittsfläche 25 cm$^2$. Wie hoch ist die Geschwindigkeit des Fluids hier?

\end{enumerate}


%\antwort{
%Kontinuitätsgleichung
%\begin{align*}
%A_1 \cdot v_1 = A_2 \cdot v_2 \\
%v_2 = \frac{A_1 \cdot v_1}{A_2}
%\end{align*}
%Einsetzen a) und b):
%\begin{align*}
%a) &&v_2 &= \frac{20\cm^2 \cdot 3\mskip3mu \cm/\text{s}}{10\mskip3mu \cm^2} &&= 6\cm/\text{s} \\
%b) &&v_2 &= \frac{20\cm^2 \cdot 3\mskip3mu \cm/\text{s}}{25\mskip3mu \cm^2} &&=2,4\cm/\text{s}
%\end{align*}
%}
%
%
%\pagebreak



%\pagebreak
%\newgeometry{top=3cm,bottom=3cm,left=4cm,right=4cm}
\subsection*{Aufgabe 02}

In einer Wasserstrahl-Vakuumpumpe, die wie ein Venturi-Rohr aufgebaut ist, fließen 4 L/min Wasser. Der Durchmesser der Zuleitung (und Ableitung) beträgt \mbox{14 mm}, der Durchmesser an der engsten Stelle 3 mm. Der absolute Druck am Ausgang der Wasserstrahl-Vakuumpumpe entspricht einem Umgebungsdruck von 1 bar. Der geodätische Höhenunterschied kann vernachlässigt werden. 
Welcher Unterdruck wird an der engsten Stelle erreicht? Welcher absolute Druck herrscht an der engsten Stelle?

\begin{marginfigure}[-3cm]
\incfig{venturi}{.9}
\caption{Beim Venturi-Rohr lässt sich die Fließgeschwindigkeit über den Höhenunterschied der angesaugten Flüssigkeit in einem Rohr ablesen, welches die Verengung mit dem normalen Rohrdurchmesser verbindet.}
\end{marginfigure}
%
%\begin{align*}
%\dot{V} &= v_2 \cdot A_2 \qquad \rightarrow v_2 = \frac{\dot{V}}{A_2} \\
%&\qquad \vdots \\
%v_2 &\approx 0,43\mskip3mu \m/\text{s}\\ 
%v_1 &= \frac{v_2 \cdot A_2}{A_1}\\
%&= v_2 \cdot \frac{d_2^2}{d_1^2} \\
%&= 0,43\mskip3mu \text{m/s} \cdot \frac{14^2}{3^2} \\
%&=9,36\mskip3mu \text{m/s} \\ \\
%\end{align*}
%Bernoulli:
%\begin{equation}
%p_1 + \frac{\rho \cdot v_1^2}{2} = p_2 + \frac{\rho \cdot v_2^2}{2}
%\end{equation}
%\begin{align*}
%p_2 - p_1 = \frac{\rho \cdot v_1^2}{2} -  \frac{\rho \cdot v_2^2}{2} 
%\end{align*}
%
%
%Werte einsetzen: 
%\begin{align*}
%p_2 - p_1 &= \frac{1000 \mskip3mu \text{kg/m}^3 \cdot (9,36\mskip3mu \text{m/s} )^2}{2} -  \frac{1000 \mskip3mu \text{kg/m}^3 \cdot (0,43 \mskip3mu \text{m/s} )^2}{2} \\
%&= 43.712 \mskip3mu \frac{\text{kg}}{\m \cdot \s^2} \\
%&= 43,7 \mskip3mu \text{kPa} \\ \\
%p_1 &= p_2 - \Delta p \\
%	&= (100 - 43,7)\text{kPa}\\ &= 56,3\mskip3mu \text{kPa}
%\end{align*}
%\pp
%
%Wobei $p_1$ der Druck an der engsten Stelle und $p_2$ der Druck am Ende des Rohres ist.\par
%Absoluter Druck $p_1$ = 56,34 \text{kPa} \par 
%Unterdruck $\Delta p$ = 43,7 \text{kPa} \par 



\subsection*{Aufgabe 3}
Wie viel Zeit braucht die Entleerung eines Rührkessels unter der Wirkung der \mbox{Erdschwere}, wenn der Bodenablass die NW 80mm aufweist? Der Rührkessel ist zu Beginn mit \mbox{10.000 L} Wasser gefüllt. Der Innendurchmesser des Rührkessels beträgt 2,4 m. Der gewölbte Boden des Klöppers ist bei der Berechnung zu vernachlässigen, ebenso der Druckverlust beim Einströmen in die Rohrleitung.






\pagebreak
%\newgeometry{top=3cm,bottom=3cm,left=4cm,right=4cm}
\subsection*{Aufgabe 4}
Wie lange dauert es, bis Wasser aus einem Gefäß mit folgenden Eigenschaften vollständig abgeflossen ist?

\begin{itemize}
\item Situation A
\begin{itemize}
    \item  Querschnittsfläche des Gefäßes: 40 cm$^2$
    \item  Querschnittsfläche der Ausflussöffnung: 2 cm$^2$
    \item  Höhe des Flüssigkeitsspiegels über dem Ausfluss: 1 m
\end{itemize}
\item Situation B
\begin{itemize}
    \item  Querschnittsfläche des Gefäßes: 5 cm$^2$
    \item  Querschnittsfläche der Ausflussöffnung: 2 cm$^2$
    \item  Höhe des Flüssigkeitsspiegels über dem Ausfluss: 1 m
\end{itemize}
\item Situation C
\begin{itemize}
    \item  Querschnittsfläche des Gefäßes: 40 cm$^2$
    \item  Querschnittsfläche der Ausflussöffnung: 2 cm$^2$
    \item  Höhe des Flüssigkeitsspiegels über dem Ausfluss: 4 m
\end{itemize}
\item Situation D
\begin{itemize}
    \item  Querschnittsfläche des Gefäßes: 40 cm$^2$
    \item  Querschnittsfläche der Ausflussöffnung: 10 cm$^2$
    \item  Höhe des Flüssigkeitsspiegels über dem Ausfluss: 1 m
\end{itemize}
\end{itemize}

\begin{marginfigure}
%\begin{figure}[ht]
\centering
\incfig{fass}{.9}
\caption{Ein Fass}
%\end{figure}
\end{marginfigure}
%\ppbig
%
%
%Torricelli:
%\begin{equation}
%t = \sqrt{\frac{2}{g}} \cdot \frac{A_1}{A_2} \cdot (\sqrt{h_a} - \sqrt{h_w})
%\end{equation}
%Für eine komplette Entleerung ($h_w$ = 0) :
%\begin{align*}
%t = \sqrt{\frac{2}{g}} \cdot \frac{A_1}{A_2} \cdot \sqrt{h_a}
%\end{align*}
%Einsetzen (Aufgabenteil a):
%\begin{align*}
%t &= 0,45 \frac{s}{\sqrt{m}}\cdot \frac{40cm^2}{2cm^2} \cdot \sqrt{1m} \\
% &= 9,03s 
%\end{align*}
%b) 1,13s ; c) 18,08s ; d) 1,81s \par 
%
%\pagebreak


%\pagebreak

%\ppbig
%
%
%\begin{align*}
%&A_1 &&= \pi \cdot (\frac{2,4\mskip3mu \m}{2})^2 &&= 4,52\mskip3mu \m^2 \\
%&A_2 &&= \pi \cdot (\frac{0,08\mskip3mu \m}{2})^2 &&= 0,005\mskip3mu \m^2 \\
%&h_a &&= \frac{V}{A_1} = \frac{10\mskip3mu \m^3}{4,52m^2} &&= 2,21\mskip3mu \m^2 
%\end{align*}
%
%\ppbig
%Torricelli:
%\begin{equation}
%t = \sqrt{\frac{2}{g}} \cdot \frac{A_1}{A_2} \cdot (\sqrt{h_a} - \sqrt{h_w})
%\end{equation}
%
%Einsetzen:
%\begin{align*}
%t &= \sqrt{\frac{2}{g}} \cdot \frac{4,52\mskip3mu \m^2}{0,005\mskip3mu \m^2} \cdot \sqrt{2,21\mskip3mu \m^2} \\
%&= 604s
%\end{align*}
%
%
%
%
%
%
%
%\pagebreak




%\pagebreak
\subsection*{Aufgabe 5}
Berechnen Sie die Reynoldszahl für Wasser bei 0°C. Ist die Strömung jeweils laminar oder turbulent?
\begin{itemize}
       \item[a.] Geschwindigkeit 5 cm/s, Rohrinnendurchmesser 30 cm
       \item[b.] Geschwindigkeit 1 cm/s, Rohrinnendurchmesser 30 cm
       \item[c.] Geschwindigkeit 5 cm/s, Rohrinnendurchmesser 10 cm
\end{itemize}



%\ppbig
%
%
%
%Kritische Reynoldszahl für turbulente oder laminare Strömung:
%\begin{equation}
%\text{Re}_{\text{krit}} = \frac{\rho \cdot d \cdot v}{\eta} = 2300
%\end{equation}
%
%
%Einsetzen:
%\begin{align*}
%&\text{a)} && \frac{1.000\mskip3mu \kg/\m^3 \cdot 0,3\mskip3mu \m \cdot 0,05\mskip3mu m/\s}{1.790 \cdot 10^{-6}\mskip3mu \text{P} \cdot \s} &&\approx 8379,9 \\
%&\text{b)} && \frac{1.000\mskip3mu \kg/\m^3 \cdot 0,3\mskip3mu \m \cdot 0,01\mskip3mu m/\s}{1.790 \cdot 10^{-6}\mskip3mu \text{P} \cdot \s} &&\approx 1676 \\
%&\text{c)} && \frac{1.000\mskip3mu \kg/\m^3 \cdot 0,1\mskip3mu \m \cdot 0,05\mskip3mu m/\s}{1.790 \cdot 10^{-6}\mskip3mu \text{P} \cdot \s} &&\approx 2793,3 \\
%\end{align*}
%
%Oberhalb der kritischen Reynolds-Zahl 2.300 erfolgt der Übergang zu turbulenten Strömungen
%
%
%




	
\end{document}
